\section{Introduction}\label{sec:introduction}

Ramsey's theorem guarantees that every sufficiently large graph contains a
large clique or independent set.  Its classical constructive proof
\cite{ErdosSzekeres1935} repeatedly chooses a pivot and restricts the
candidate set to one of its two colour neighbourhoods.  Starting from
$M=4^{K-1}$ vertices, truncating each retained neighbourhood to half its
current size gives an $O(M)$-query algorithm.

The central observation of this paper is that the candidate sets need not be
materialized.  After $i$ rounds, membership in the current set is determined
by the $i$ adjacency constraints imposed by the previous pivots.  These short
predicates allow quantum search to access the same recursion at substantially
smaller query cost.  The threshold $4^{K-1}$ is the natural scale for this
recursive construction.

\subsection{Our results}

We work in the promised valid-graph adjacency-oracle model.  One coherent
query returns the colour of a requested unordered vertex pair, and the output
is a set of vertices inducing one colour.  Our main result is the following.

\begin{theorem}[Implicit-majority quantum Ramsey search]
\label{thm:main}
Let $K\ge2$, $N\ge4^{K-1}$, and $0<\eta<1/2$.  Given coherent edge-oracle
access to a simple undirected graph on $N$ vertices, there is a quantum
algorithm which, with probability at least $1-\eta$, outputs a $K$-clique or
a $K$-vertex independent set using
\[
 O\!\left(2^K K\log\frac K\eta\right)
\]
edge queries in the worst case.  The algorithm verifies every returned
witness and may output $\bot$ on the exceptional event.
\end{theorem}

At the succinct Ramsey scale, the theorem gives the following form.

\begin{corollary}\label{cor:succinct}
Let $n\ge2$, $N=2^n$, and $0<\eta<1/2$.  Given coherent edge-oracle access
to a simple graph on $N$ vertices, one can find a clique or independent set
of order $\lfloor n/2\rfloor+1$ using
\[
 O\!\left(
 \sqrt N\log N\log\frac{\log N}{\eta}
 \right)
\]
edge queries in the worst case.
\end{corollary}

\begin{proof}
Set $K=\lfloor n/2\rfloor+1$.  Then
\[
 4^{K-1}=2^{2\lfloor n/2\rfloor}\le 2^n=N,
\]
so \cref{thm:main} applies.  Moreover, $2^K\le2\sqrt N$, $K=O(\log N)$, and
$\log(K/\eta)=O(\log(\log N/\eta))$.  Substituting these estimates into the bound of
\cref{thm:main} proves the claim.
\end{proof}

We also obtain an estimation-free quantum algorithm with query complexity
$O(2^K K^2\log K\log(1/\eta))$, a multicolour extension of that algorithm,
and a randomized classical lower bound of
$\Omega(M^{1-1/\sqrt2})$ when $M=4^{K-1}$.  These results appear in
\cref{prop:size-biased-algorithm,cor:multicolour,prop:classical-lower}.

\subsection{Technical overview}

Both algorithms represent the candidate set after round $i$ by the
conjunction of the adjacency constraints generated by the first $i$ pivots.
A membership test therefore uses $O(i)$ edge queries.  Capped
unknown-solution search samples from such a set without first enumerating its
vertices; conditional on success, permutation symmetry makes the sample
exactly uniform.  The precise sampler guarantee and its proof appear in
\cref{lem:capped-search,app:capped-search}.

The first recursion removes majority estimation entirely.  At each pivot, it
samples a uniform candidate and follows the colour of the sampled edge.  A
colour class is therefore selected with probability proportional to its
size.  The survival probability obeys a one-dimensional recurrence, and
\cref{lem:size-biased-survival} shows that the recursion completes with
probability greater than $1/2$ from $4^{K-1}$ vertices.  This yields the
estimation-free bound in \cref{prop:size-biased-algorithm}.

The sharper recursion estimates a majority colour at each pivot.  Sampling a
deep candidate set is exponentially more expensive than sampling an early
one, so imposing the same estimation accuracy at every level wastes queries.
We instead distribute a fixed total error budget across the recursion.  Early
majorities receive higher accuracy, while later majorities tolerate larger
error.  A dyadic schedule keeps the cumulative loss in candidate-set size
bounded and balances the work across levels.  Conditional concentration,
the resulting density invariant, and exact output correctness are proved in
\cref{lem:concentration,lem:survival,lem:output}; the query summation in
\cref{sec:complexity} then proves \cref{thm:main}.

\subsection{Comparison with a reported lower bound}

At the default $N=2^n$, $K=n/2$ parameters,
\cite[Section III, p.~415]{JainLiRobereXun2024} reports an
$N^{1-o(1)}$ quantum-query lower bound for the same Ramsey relation.  This
creates an apparent discrepancy with \cref{cor:succinct}.
Section~\ref{sec:prior-art} aligns the two formulations and shows that direct
substitution into the cited ingredients yields an exponent at most $1/24$,
rather than $1-o(1)$.

\paragraph{Organization.}
\Cref{sec:model} defines the query model and sampler, and
\cref{sec:algorithm,sec:correctness,sec:complexity} present the algorithms and
proofs.  \Cref{sec:prior-art} compares the result with prior work, and
\cref{sec:scope} concludes with open directions.  The appendices contain the
sampler proof, the classical lower-bound derivation, the multicolour proof,
and numerical illustrations.
